\documentclass[11pt]{article}
\usepackage{amsmath,amssymb,hyperref}
\title{A requested radio model is not evidence that a cave test ran it}
\author{}\date{}
\begin{document}\maketitle
\paragraph{Source cut.} The pinned SubT
\href{https://github.com/osrf/subt/blob/083c6a8acb0c25aa1a35e34517e628cd3b52ae77/subt_ign/launch/cloudsim_sim.ign}
{launch template} requests visibility-range radio. The exact broker
selection installs that model only if its visibility model initialized.
The constructor depends on reading a visibility graph and its sampled
lookup table; an unavailable table permits the selected range-only default
to remain in place. This is a startup distinction. It does not imply
per-send fallback in an already initialized visibility model. In particular
the template's \texttt{generate\_table} setting does not establish that a
missing table will be generated here: the generation branch in the pinned
broker loader is commented out.

\paragraph{Matched boundary.} Replay the selected original C++ loader
operation on a staged valid table and on an absent file, then run the
brokers' chosen original power callbacks at identical endpoint poses.
Hold the request, radio constants and world declarations fixed. Exercise
the native sampled-point lookup rather than assuming poses identify exact
visibility vertices. The actual \texttt{cave\_circuit\_01} DOT and SDF
supply the weighted graph and model origins. A small, explicitly staged
lookup table associates unambiguous rounded model origins with their DOT
tile IDs. This deliberately does not claim to have the stock binary LUT,
the actual mesh point classification, full Gazebo, or recorded corridor
motion. Static graph costs are computed separately from the pinned DOT;
an independent direct arithmetic model checks both power functions.
The construction tests a configuration dependency, not how to evade any
operational control.

For a mapped no-relay link, let \(c\) be the stored graph cost and \(d\)
the direct distance. With both poses the same in the two runs, the
range-only mode gives finite mean power for \(0<d\le500\). Under the
visibility model without a relay it also needs a known cost \(c\le10\),
and otherwise returns minus infinity. The base to the chosen sampled
tile 2 is \(76.576759\) units away with graph cost \(16\): finite range mean
\(-48.261455\) versus minus infinity. A mapped cost-1 control is finite in
both modes. An unmapped pose \(92\) units from the base also differs,
and a separate unmapped \(592\)-unit case is rejected by both. Across
24,182 ordered no-relay queries, including those two controls, 1,046
were finite under both, 15,513 finite under range only, and 7,623 finite
under neither. These are counts for the staged test cases, not population
estimates about links in the cave.

The direction can reverse once a route through relay nodes is given.
The separate staged relay row uses edges from the same actual DOT and
origins from the same SDF: direct distance \(500.388849\), max hop
\(33\), 28 staged relays and max graph hop cost \(6\). The visibility
callback returns finite mean \(-72.633530\); range-only rejects the direct
path. This is a conditional test of the native visibility RF consumer,
not a claim that the stock chooser selects that route or that the
benchmark supplies 28 deployable relays. Consequently the mode replacement
is not, in general, a containment in either direction for finite-mean
eligibility.

\paragraph{Consumption and limits.} The relevant outcome here is whether
the selected RF function deterministically excludes a candidate link at
the power layer. Finite mean is not an application receipt, end-to-end
delivery probability or deadline guarantee. The fallback function does
not populate the optional relay-use output of the tested callback binding;
the two initialized output values in the unit replay stay different.
The replay calls the power functions with initialized output storage.
It does not exercise packet dispatch or infer a runtime value for an
out-parameter not set by that function.

A configuration-only attestation cannot distinguish these two runs.
A replay intended to exercise disabled contact must record the effective
selected model and which table was actually loaded, then check the
intended absence/presence of contact and the application outcome on the
relevant trace. A read-only control comparing to the same effective
model sees the same behavior. A dynamic obstacle can also change
travel while this static graph/RF substrate has not changed; no result
here infers that a rockfall blocks mobility or radio, or that an obstacle
asset cannot change the live poses or the relay inputs. This is a held-in
I2 evidence gate, not an independent mission or a coordination gain.

\end{document}
